\section{Conclusioni e sviluppi futuri}
\label{sec:conclusioni}
\thispagestyle{plain}

L'architettura progettata dimostra concretamente la fattibilità e l'efficacia dell'integrazione di strumenti open-source eterogenei per l'implementazione di un ecosistema Zero Trust maturo. Il presente lavoro di tesi ha permesso di superare il tradizionale modello di sicurezza perimetrale, realizzando un'infrastruttura in cui la fiducia non è mai presunta, bensì verificata e quantificata continuamente in tempo reale.

\subsection{Obiettivi raggiunti}
Il sistema ha soddisfatto appieno i requisiti del paradigma Zero Trust, declinandoli attraverso tre pilastri architetturali fondamentali:

\begin{itemize}
    \item \textbf{Segregazione e Sicurezza Infrastrutturale (PEP):} La rete è stata blindata operando una rigorosa microsegmentazione L3/L4 tramite \texttt{nftables}, forzando il transito delle comunicazioni esclusivamente attraverso il proxy \texttt{envoy}. A livello di trasporto, l'imposizione della mutua autenticazione crittografica (\textit{mTLS}) e l'estrazione dell'impronta digitale \textit{JA3} hanno permesso di verificare con certezza crittografica l'identità del client e la conformità dell'hardware sottostante (tramite attestazione TPM).
    
    \item \textbf{Autorizzazione Continua e Difesa in Profondità (PDP):} Il motore decisionale \texttt{open policy agent} (OPA) ha centralizzato l'enforcement delle policy valutando dinamicamente le sei dimensioni contestuali (ZTA 6D). L'implementazione di clausole di guardia (\textit{Short-Circuiting}) ha ottimizzato i tempi di risposta bloccando istantaneamente gli accessi non conformi. Inoltre, l'applicazione del principio di \textit{Defense in Depth} ha garantito la protezione del dato sensibile attraverso due livelli di ispezione profonda (L7 DPI): un filtro \texttt{Lua} per la mitigazione rapida degli attacchi HTTP e l'ispezione nativa dei payload BSON di \texttt{mongodb} per la neutralizzazione di query distruttive.
    
    \item \textbf{Intelligenza Predittiva e Auditing (SIEM):} L'integrazione con lo \texttt{splunk} Machine Learning Toolkit (MLTK) ha segnato il definitivo passaggio da un'autenticazione statica a una \textit{Continuous Authorization} comportamentale. L'impiego dell'algoritmo \textit{Gradient Boosting} ha permesso di correlare 12 feature contestuali e comportamentali, restituendo un rischio predittivo accurato in grado di smascherare minacce asimmetriche e \textit{Insider Threat}, adattando le soglie di tolleranza alla reale postura di sicurezza del dispositivo. Contemporaneamente, le sonde \texttt{snort} hanno garantito un monitoraggio continuo contro i tentativi di elusione e ricognizione.
\end{itemize}

\subsection{Sviluppi futuri}
L'infrastruttura realizzata costituisce una base solida e scalabile, aperta a diverse evoluzioni ingegneristiche in ambito Enterprise. Tra i principali sviluppi futuri si individuano:

\begin{enumerate}
    \item \textbf{Gestione dinamica delle Identità (SPIFFE/SPIRE):} L'attuale generazione statica dei certificati x509 potrebbe essere sostituita dall'integrazione dello standard \textit{SPIFFE/SPIRE}. Questo consentirebbe il rilascio di certificati a brevissima scadenza (\textit{short-lived certificates}) basati sull'attestazione continua del workload, eliminando la complessità operativa legata alla gestione e distribuzione delle \textit{Certificate Revocation List} (CRL).
    
    \item \textbf{Integrazione di dispositivi Legacy e M2M (Kiosk):} Il modello ZTA attuale prevede l'obbligo tassativo del modulo TPM per l'autenticazione. Uno sviluppo futuro prevede l'estensione sicura della rete a dispositivi operativi sprovvisti di hardware crittografico dedicato (es. chioschi informativi o comunicazioni \textit{Machine-to-Machine} via API). Ciò richiederà l'addestramento di modelli ML specifici che compensino l'assenza del TPM abbassando drasticamente le soglie di tolleranza e monitorando con maggiore severità l'ancoraggio all'impronta \textit{JA3}.
    
    \item \textbf{Sensor Fusion per il Machine Learning (Integrazione L3/L4):} Attualmente, il \textit{Trust Score} di Splunk MLTK è calcolato su 12 feature derivanti dal livello applicativo (Envoy e Web API). Un potenziamento cruciale consisterà nell'ingestione della telemetria di rete a basso livello all'interno del vettore di addestramento. Trasformare gli allarmi di \texttt{snort} (es. tentativi di scansione Nmap) e i drop log di \texttt{nftables} in feature comportamentali addizionali permetterebbe al modello \textit{Gradient Boosting} di correlare anomalie applicative con movimenti laterali di rete, aumentando esponenzialmente la precisione nel rilevamento di minacce persistenti avanzate (APT).

    \item \textbf{Orchestrazione e Risposta Automatica (SOAR):} L'architettura prevede già un'interfaccia API REST (\texttt{FastAPI}) per l'inserimento dinamico di IP malevoli nella \textit{denylist} di \texttt{nftables}, attualmente utilizzabile dal team SOC per mitigazioni manuali. Questa componente è progettata per evolvere in una pipeline SOAR (\textit{Security Orchestration, Automation, and Response}) completa: istruendo Splunk a lanciare \textit{webhook} automatici verso l'API del firewall non appena il rischio comportamentale supera una soglia critica, si garantirà l'isolamento istantaneo della minaccia a livello kernel (L3) senza alcun intervento umano.
    
    \item \textbf{Federated Machine Learning:} In scenari multi-ospedale, il modello \textit{Gradient Boosting} potrebbe evolvere verso paradigmi di \textit{Federated Learning}, permettendo ai SIEM di diverse sedi ospedaliere di addestrare un modello globale condividendo esclusivamente i gradienti di errore e i pesi algoritmici, senza mai scambiarsi i log grezzi contenenti dati clinici e rispettando in modo nativo i vincoli normativi sulla privacy (GDPR).
\end{enumerate}